\section{Summary}

\begin{frame}{What to Take Away}

    \begin{itemize}
        \setlength\itemsep{0.1em}
        \item \textbf{Aggregate metrics hide subgroup failure.} The single most valuable habit from
              this lecture is the disaggregated evaluation table, with sample sizes, including
              intersections.

        \item \textbf{Bias is mostly a specification problem.} It enters through the target
              definition, the label, the sampling frame, and the deployment context --- not only
              through ``the data.'' Dropping the sensitive attribute does not remove it.

        \item \textbf{The fairness definitions genuinely conflict.} When base rates differ,
              calibration and error-rate balance cannot both hold. Choose the criterion that
              matches the harm, and write down what you gave up.

        \item \textbf{Explanations describe the model, not the world.} SHAP and LIME are excellent
              for debugging, sensitive to correlated features and to the background distribution,
              and not causal. Validate a hypothesis before acting on it.

        \item \textbf{Documentation is the deliverable.} The model card and the audit are what a
              reviewer, a customer, or a regulator will actually read --- and under the EU AI Act,
              much of their content is now required for high-risk systems.
    \end{itemize}
\end{frame}

\begin{frame}{A Checklist You Can Use on Monday}

    \begin{columns}[T]
        \begin{column}{0.49\textwidth}
            \textbf{Before training}
            \begin{enumerate}\scriptsize
                \setlength\itemsep{0.28em}
                \item Write one sentence: who is affected, and how could they be harmed?
                \item Ask what your label \emph{actually} measures, and whether that differs by
                      group.
                \item Check the composition of your data against the intended use population.
                \item Fit an interpretable baseline and record its score.
            \end{enumerate}

            \vspace{0.6em}
            \textbf{Before shipping}
            \begin{enumerate}\scriptsize
                \setlength\itemsep{0.28em}
                \setcounter{enumi}{4}
                \item Produce the disaggregated table with $n$ and intervals.
                \item Name the fairness criterion and what you traded away.
                \item Run SHAP or permutation importance and look for a shortcut feature.
            \end{enumerate}
        \end{column}
        \begin{column}{0.49\textwidth}
            \textbf{At release}
            \begin{enumerate}\scriptsize
                \setlength\itemsep{0.28em}
                \setcounter{enumi}{7}
                \item Write the model card, including out-of-scope uses and the worst slice.
                \item Get a review from someone who did not build it.
                \item Set the re-audit date and the monitoring alert on disparity.
            \end{enumerate}

            \vspace{1.0em}
            \begin{tcolorbox}[colback=raiGreen!5, colframe=raiGreen!60, boxrule=0.7pt]
                \scriptsize
                Items 1, 2 and 5 catch most of what goes wrong, and none of them require a new
                library.
            \end{tcolorbox}
        \end{column}
    \end{columns}

    \vspace{0.8em}
    \centering
    \scriptsize \textbf{Hands-on:} the SHAP/LIME and fairness-audit exercises described on the
    preceding slides (companion notebooks in preparation).
\end{frame}
